% GENERATED FILE. Edit canonical lessons, then run npm run generate:books.
% canonical-source-hash: fnv1a64:c5da7684952108eb
% canonical-lessons: LA-C400-doceo, LA-C400-erudio, LA-C400-instituo, LA-C400-instruo, LA-C400-praecipio

\chapter{Teaching}
\label{ch:la-teaching}
\hlchaptermodality{400}

\begin{chapteropening}
\textbf{By the end of this chapter:} \emph{I can say that someone teaches, educates, trains and instructs a pupil, and tells him what to do.}

Complete the last lesson of chapter 400: I can say that someone teaches, educates, trains and instructs a pupil, and tells him what to do.
\end{chapteropening}

\section[doceō, docēre]{doceō, docēre — to teach}

\label{lesson:LA-C400-doceo}



\begin{quote}
Before the new one: say the Latin for to deceive, then the Latin for to call to witness.
\end{quote}

\subsection*{You'll want to know: doceō, docēre}
\textbf{doceō, docēre} — \textquotedblleft{}to teach\textquotedblright{}.

\begin{quote}
Magister puerōs litterās docet.
\end{quote}

\begin{quote}
\emph{The teacher teaches the boys their letters.}
\end{quote}

\emph{doceō, docēre, docuī}; its participle follows in a later lesson. Two accusatives, the person and the thing: \emph{puerōs litterās docet}. English \emph{doctor}, first a teacher, and \emph{docile}.

\begin{cousinweb}
\emph{doceō} is inherited from Proto-Italic \emph{*dokeō} (EtymDB), usually taken as a causative beside \emph{decet} \textquotedblleft{}it is fitting\textquotedblright{}: to make someone take a thing in. English \emph{doctor} and \emph{doctrine} come from it.
\end{cousinweb}

\begin{grammarlens}[title={What you've built}]
One more word for telling, explaining and describing at B1.
\end{grammarlens}

\subsection*{Your turn}
\begin{itemize}
\raggedright
  \item \emph{Say it:} \emph{doceō, docēre}
  \item \emph{Say it:} \emph{doceō, docēre}, once more
  \item \emph{Say it:} a sentence of your own with \emph{doceō, docēre}
  \item \emph{Recall:} say the Latin for to deceive, then the Latin for to call to witness, then say \emph{doceō, docēre} again
\end{itemize}

\begin{tcolorbox}[breakable,colback=teal!4,colframe=teal!35!black,title={Before you move on}]
What does doceō, docēre mean? (\textquotedblleft{}To teach\textquotedblright{}.) Say it once more.
\end{tcolorbox}
\section[ērudiō, ērudīre]{ērudiō, ērudīre — to educate, to polish}

\label{lesson:LA-C400-erudio}



\begin{quote}
Before the new one: say the Latin for to call to witness, then the Latin for to teach.
\end{quote}

\subsection*{You'll want to know: ērudiō, ērudīre}
\textbf{ērudiō, ērudīre} — \textquotedblleft{}to educate, to polish\textquotedblright{}.

\begin{quote}
Pater fīlium dīligenter ērudīvit.
\end{quote}

\begin{quote}
\emph{The father educated his son carefully.}
\end{quote}

\emph{ērudiō, ērudīre, ērudīvī}; its participle follows in a later lesson. \emph{Ē} + \emph{rudis}, rough: to take the roughness out. English \emph{erudite}.

\begin{cousinweb}
\emph{ērudiō} is built from \emph{ē} \textquotedblleft{}out of\textquotedblright{} and \emph{rudis} \textquotedblleft{}rough, untaught\textquotedblright{} (EtymDB; Lewis \& Short): to bring out of the rough state. English \emph{erudite} comes from its participle.
\end{cousinweb}

\begin{grammarlens}[title={What you've built}]
One more word for telling, explaining and describing at B1.
\end{grammarlens}

\subsection*{Your turn}
\begin{itemize}
\raggedright
  \item \emph{Say it:} \emph{ērudiō, ērudīre}
  \item \emph{Say it:} \emph{ērudiō, ērudīre}, once more
  \item \emph{Say it:} a sentence of your own with \emph{ērudiō, ērudīre}
  \item \emph{Recall:} say the Latin for to call to witness, then the Latin for to teach, then say \emph{ērudiō, ērudīre} again
\end{itemize}

\begin{tcolorbox}[breakable,colback=teal!4,colframe=teal!35!black,title={Before you move on}]
What does ērudiō, ērudīre mean? (\textquotedblleft{}To educate, to polish\textquotedblright{}.) Say it once more.
\end{tcolorbox}
\section[īnstituō, īnstituere]{īnstituō, īnstituere — to train up, to instruct; to set up}

\label{lesson:LA-C400-instituo}



\begin{quote}
Before the new one: say the Latin for to teach, then the Latin for to educate.
\end{quote}

\subsection*{You'll want to know: īnstituō, īnstituere}
\textbf{īnstituō, īnstituere} — \textquotedblleft{}to train up, to instruct; to set up\textquotedblright{}.

\begin{quote}
Puerōs ad virtūtem īnstituēbat.
\end{quote}

\begin{quote}
\emph{He trained the boys in virtue.}
\end{quote}

\emph{īnstituō, īnstituere, īnstituī, īnstitūtum}, \emph{in} + \emph{statuō}: to set up, then to set a young mind in order. \emph{Īnstitūtum} is its noun.

\begin{cousinweb}
\emph{īnstituō} is built from \emph{in} and \emph{statuō} \textquotedblleft{}to set, to place\textquotedblright{} (EtymDB; Lewis \& Short), the vowel weakened inside the compound. English \emph{institute} comes from it.
\end{cousinweb}

\begin{grammarlens}[title={What you've built}]
One more word for telling, explaining and describing at B1.
\end{grammarlens}

\subsection*{Your turn}
\begin{itemize}
\raggedright
  \item \emph{Say it:} \emph{īnstituō, īnstituere}
  \item \emph{Say it:} \emph{īnstituō, īnstituere}, once more
  \item \emph{Say it:} a sentence of your own with \emph{īnstituō, īnstituere}
  \item \emph{Recall:} say the Latin for to teach, then the Latin for to educate, then say \emph{īnstituō, īnstituere} again
\end{itemize}

\begin{tcolorbox}[breakable,colback=teal!4,colframe=teal!35!black,title={Before you move on}]
What does īnstituō, īnstituere mean? (\textquotedblleft{}To train up, to instruct; to set up\textquotedblright{}.) Say it once more.
\end{tcolorbox}
\section[īnstruō, īnstruere]{īnstruō, īnstruere — to equip, to draw up; to instruct}

\label{lesson:LA-C400-instruo}



\begin{quote}
Before the new one: say the Latin for to educate, then the Latin for to train up.
\end{quote}

\subsection*{You'll want to know: īnstruō, īnstruere}
\textbf{īnstruō, īnstruere} — \textquotedblleft{}to equip, to draw up; to instruct\textquotedblright{}.

\begin{quote}
Imperātor aciem īnstrūxit.
\end{quote}

\begin{quote}
\emph{The general drew up his battle line.}
\end{quote}

\emph{īnstruō, īnstruere, īnstrūxī, īnstrūctum}: to build in, to equip, then to fit someone out with knowledge. English \emph{instruct} and \emph{instrument}.

\begin{cousinweb}
\emph{īnstruō} is built from \emph{in} and \emph{struō} \textquotedblleft{}to pile up, to build\textquotedblright{} (Lewis \& Short). English \emph{instruct} and \emph{instrument} come from it.
\end{cousinweb}

\begin{grammarlens}[title={What you've built}]
One more word for telling, explaining and describing at B1.
\end{grammarlens}

\subsection*{Your turn}
\begin{itemize}
\raggedright
  \item \emph{Say it:} \emph{īnstruō, īnstruere}
  \item \emph{Say it:} \emph{īnstruō, īnstruere}, once more
  \item \emph{Say it:} a sentence of your own with \emph{īnstruō, īnstruere}
  \item \emph{Recall:} say the Latin for to educate, then the Latin for to train up, then say \emph{īnstruō, īnstruere} again
\end{itemize}

\begin{tcolorbox}[breakable,colback=teal!4,colframe=teal!35!black,title={Before you move on}]
What does īnstruō, īnstruere mean? (\textquotedblleft{}To equip, to draw up; to instruct\textquotedblright{}.) Say it once more.
\end{tcolorbox}
\section[praecipiō, praecipere]{praecipiō, praecipere — to instruct, to advise; to take beforehand}

\label{lesson:LA-C400-praecipio}



\begin{quote}
Before the new one: say the Latin for to train up, then the Latin for to equip.
\end{quote}

\subsection*{You'll want to know: praecipiō, praecipere}
\textbf{praecipiō, praecipere} — \textquotedblleft{}to instruct, to advise; to take beforehand\textquotedblright{}.

\begin{quote}
Medicus mihi praecēpit ut quiēscerem.
\end{quote}

\begin{quote}
\emph{The doctor told me to rest.}
\end{quote}

\emph{praecipiō, praecipere, praecēpī}; \emph{prae} + \emph{capiō}: to take beforehand, then to lay down a rule: \emph{praecipiō tibi ut}, I instruct you to. English \emph{precept}.

\begin{cousinweb}
\emph{praecipiō} is built from \emph{prae} \textquotedblleft{}before\textquotedblright{} and \emph{capiō} \textquotedblleft{}to take\textquotedblright{} (EtymDB; Lewis \& Short): to take in advance, then to tell in advance. English \emph{precept} comes from it.
\end{cousinweb}

\begin{grammarlens}[title={What you've built}]
One more word for telling, explaining and describing at B1.
\end{grammarlens}

\subsection*{Your turn}
\begin{itemize}
\raggedright
  \item \emph{Say it:} \emph{praecipiō, praecipere}
  \item \emph{Say it:} \emph{praecipiō, praecipere}, once more
  \item \emph{Say it:} a sentence of your own with \emph{praecipiō, praecipere}
  \item \emph{Recall:} say the Latin for to train up, then the Latin for to equip, then say \emph{praecipiō, praecipere} again
\end{itemize}

\begin{tcolorbox}[breakable,colback=teal!4,colframe=teal!35!black,title={Before you move on}]
What does praecipiō, praecipere mean? (\textquotedblleft{}To instruct, to advise; to take beforehand\textquotedblright{}.) Say it once more.
\end{tcolorbox}
